% Electricity Lab — лабораторный практикум.
% Сборка: tectonic electricity-lab-manual.tex   (движок XeTeX, кириллица + системные DejaVu)
% Источник параллелен electricity-lab-manual.html; картинки — из img/ (режим --capture).
\documentclass[11pt]{article}

\usepackage{fontspec}
\usepackage{geometry}
\geometry{a4paper, top=24mm, bottom=22mm, left=20mm, right=20mm}

% Системные шрифты (стоят в Windows\Fonts), те же, что в CSS-версии.
\setmainfont{DejaVu Sans}[Scale=0.92]
\newfontfamily\seriffont{DejaVu Serif}[Scale=0.92]

% babel (а не polyglossia): под XeLaTeX даёт русские переносы и НЕ требует
% OpenType-тег «Cyrl» в каждом шрифте (DejaVu Serif имеет глифы, но не тег).
\usepackage[english,russian]{babel}

\usepackage{amsmath}
\usepackage{amssymb}
\usepackage{graphicx}
\usepackage{xcolor}
\usepackage[most]{tcolorbox}
\usepackage{booktabs}
\usepackage{array}
\usepackage{enumitem}
\usepackage{titlesec}
\usepackage{microtype}
\usepackage[hidelinks]{hyperref}

% ---- палитра (из CSS) ----
\definecolor{accent}{HTML}{2B6CB0}
\definecolor{ink}{HTML}{1A1D24}
\definecolor{muted}{HTML}{5B6472}
\definecolor{warn}{HTML}{B0472F}
\definecolor{linecol}{HTML}{D4D9E0}
\definecolor{formulabg}{HTML}{F3F6FA}
\definecolor{predictbg}{HTML}{FFF8EC}\definecolor{predictbar}{HTML}{E0A92A}
\definecolor{observebg}{HTML}{EEF6FF}\definecolor{observebar}{HTML}{2B6CB0}
\definecolor{explainbg}{HTML}{EEFAF0}\definecolor{explainbar}{HTML}{2F9E5A}

\color{ink}
\setlength{\parindent}{0pt}
\setlength{\parskip}{5pt}
\linespread{1.12}

% ---- заголовки ----
\titleformat{\section}{\color{accent}\Large\bfseries}{}{0em}{}[{\color{linecol}\titlerule[2pt]}]
\titlespacing*{\section}{0pt}{2pt}{10pt}
% h3 — мелкий капсом
\newcommand{\hh}[1]{\par\medskip{\footnotesize\bfseries\color{ink}\MakeUppercase{#1}}\par\nobreak\vspace{2pt}}
% лид-абзац
\newcommand{\lead}[1]{{\color{muted}\large #1}\par}

% ---- бейдж-тег рядом с названием работы ----
\newtcbox{\badge}{on line, colback=accent, colframe=accent, coltext=white,
  boxrule=0pt, arc=4pt, left=5pt, right=5pt, top=1pt, bottom=1pt, fontupper=\footnotesize}
% заголовок работы: «Работа N» + тег
\newcommand{\lab}[2]{\section*{Работа #1\hspace{8pt}\badge{#2}}}

% ---- формула-пилюля (как .f в CSS) ----
\newtcbox{\f}{on line, colback=formulabg, colframe=formulabg, boxrule=0pt, arc=3pt,
  left=5pt, right=5pt, top=1pt, bottom=1pt, fontupper=\seriffont}
\newcommand{\fm}[1]{\f{\ensuremath{#1}}}     % формула в пилюле
\newcommand{\vr}[1]{\ensuremath{#1}}          % переменная курсивом (как .v)
\newcommand{\misc}[1]{\textcolor{warn}{\textbf{#1}}}

% ---- блоки «предскажи / наблюдай / объясни» ----
\tcbset{poe/.style={breakable, enhanced, frame hidden, boxrule=0pt, arc=3pt,
  left=10pt, right=10pt, top=7pt, bottom=7pt, before skip=8pt, after skip=8pt}}
\newtcolorbox{predict}{poe, colback=predictbg, borderline west={3.5pt}{0pt}{predictbar}}
\newtcolorbox{observe}{poe, colback=observebg, borderline west={3.5pt}{0pt}{observebar}}
\newtcolorbox{explain}{poe, colback=explainbg, borderline west={3.5pt}{0pt}{explainbar}}

% ---- рисунки ----
\setlength{\fboxsep}{0pt}
\newcommand{\fig}[2]{\begin{center}%
  \fcolorbox{linecol}{white}{\includegraphics[width=0.80\linewidth]{#1}}\\[3pt]%
  {\footnotesize\color{muted} #2}\end{center}}
% два рисунка в ряд
\newcommand{\figtwo}[4]{\begin{center}%
  \begin{minipage}[t]{0.485\linewidth}\centering
    \fcolorbox{linecol}{white}{\includegraphics[width=\linewidth]{#1}}\\[3pt]
    {\footnotesize\color{muted} #2}\end{minipage}\hfill
  \begin{minipage}[t]{0.485\linewidth}\centering
    \fcolorbox{linecol}{white}{\includegraphics[width=\linewidth]{#3}}\\[3pt]
    {\footnotesize\color{muted} #4}\end{minipage}\end{center}}

\graphicspath{{img/}}
\setlist{nosep, leftmargin=18pt, topsep=3pt, itemsep=3pt}

\begin{document}
\sloppy

% ================= ОБЛОЖКА =================
\thispagestyle{empty}
\vspace*{55mm}
{\color{accent}\fontsize{40}{46}\selectfont\bfseries Electricity\,Lab}\par
\vspace{6pt}
{\color{muted}\Large Лабораторный практикум по электрическим цепям}\par
\vspace{4pt}
{\color{muted}\large 10 работ \textperiodcentered{} схемы, формулы, метод «предскажи → наблюдай → объясни»}\par
\vspace{34pt}
{\color{muted}\small Иллюстрации сняты прямо из симулятора (режим \texttt{-{}-capture}).\par
Версия 2026-06-15. Решатель: MNA (узловой анализ), транзиент — companion-модели.}
\clearpage

% ================= ИНТРО =================
\section*{Как работать с практикумом}
\lead{Каждая работа открывает одну идею. Симулятор честно решает цепь (метод узловых
потенциалов), а вы проверяете свою интуицию \textbf{до} того, как увидите ответ.}

\hh{Метод «предскажи → наблюдай → объясни»}
\begin{enumerate}
  \item \textbf{Предскажи.} Прежде чем запускать симуляцию — запиши свой прогноз (ток, напряжение,
        что ярче/быстрее). Это обязательный шаг: без него мозг не учится.
  \item \textbf{Наблюдай.} Открой демо-схему в симуляторе, запусти, сними показания.
  \item \textbf{Объясни.} Совпало? Если нет — найди, какое допущение было неверным. Расхождение
        ценнее совпадения.
\end{enumerate}

\hh{Условные обозначения}
\begin{itemize}
  \item Формула в рамке: \fm{V = I\cdot R}
  \item \misc{Красным} — типичное заблуждение, по которому бьёт работа.
  \item Жёлтый блок — предсказание; синий — наблюдение; зелёный — объяснение.
\end{itemize}

\hh{Три картины одной цепи}
Симулятор показывает одну и ту же физику в нескольких метафорах. В практикуме используются:
\textbf{электрическая} схема (стрелки тока, карта потенциала) и \textbf{механическая} аналогия
(цепь-цепочка: ток = скорость звеньев, сопротивление = тормоз, конденсатор = пружина,
катушка = маховик). Аналогии помогают «пощупать» абстракцию.

\hh{Содержание}
\begin{enumerate}
  \item Закон Ома
  \item Делитель напряжения
  \item Нагруженный делитель: эффект нагрузки
  \item Мост Уитстона
  \item Зарядка конденсатора (RC, постоянная времени)
  \item Нарастание тока в катушке (RL)
  \item Колебательный контур RLC
  \item Диод: односторонняя проводимость
  \item Полуволновой выпрямитель (AC → DC)
  \item RC-фильтр нижних частот
\end{enumerate}
\clearpage

% ================= ТЕОРИЯ =================
\section*{Теория \textperiodcentered{} ток, потенциал и поле}
\lead{Краткая теория ровно по тем явлениям, которые симулятор честно считает и
показывает слоями. Формулы — рабочие; вывод через интегралы оставлен учебникам.}

\hh{Заряд и ток}
Электрический ток — упорядоченное движение заряда. Сила тока \fm{I = dq/dt} [А] — это
заряд, проходящий через сечение за секунду. В металле ток несут электроны, но за
положительное направление исторически принят \textbf{конвенциональный ток} — от «+» к
«−» во внешней цепи, навстречу движению электронов. Симулятор рисует оба: слой
\textit{current} (конвенциональный, стрелки) и \textit{electronFlow} (поток электронов,
в обратную сторону).

\hh{Дрейф электронов}
Электроны и без тока летают хаотично (тепловое движение, \ensuremath{\sim}10\textsuperscript{5}~м/с),
но в сумме никуда не смещаются. Поле добавляет крошечную направленную \textbf{дрейфовую}
скорость \fm{v_d \propto I/A} (\vr{A} — площадь сечения). Она ничтожна (доли мм/с) — но
именно она и есть ток. В слое \textit{drift} частицы дрейфуют со скоростью, пропорциональной
току (при \vr{I}=0 — стоят), а в теле резистора видна «решётка Друде»: электроны бьются о
ионы — это и есть сопротивление.

\hh{Потенциал и напряжение}
Потенциал \vr{\varphi} [В] — потенциальная энергия единичного заряда в данной точке.
Напряжение \fm{U = \varphi_1 - \varphi_2} — разность потенциалов, «толкающая» заряд.
Земля (\textit{ground}) задаёт \vr{\varphi}=0. Вдоль однородного провода потенциал падает
линейно; слой \textit{potential} красит провод в цвет потенциала.

\hh{Электрическое поле в проводнике}
Поле — это «крутизна» потенциала: \fm{E = -d\varphi/dx}. Внутри прямого участка
\fm{E = U/L} (\vr{U} — падение на длине \vr{L}). Поле направлено от высокого потенциала к
низкому и гонит заряд. Слой \textit{electricField} рисует стрелки \vr{E}, а \textit{LIC} —
сплошную текстуру силовых линий.

\hh{Поверхностный заряд}
Откуда поле внутри провода «знает», куда течь? Его создаёт тонкий слой заряда на
поверхности: где потенциал выше — избыток «+», ниже — «−». Плотность (следствие теоремы
Гаусса) \fm{\sigma \propto \varepsilon_0 E \propto \Delta U / L}. Слой \textit{surfaceCharge}
ставит на кромках провода \ensuremath{\pm} кружки.
\clearpage

\section*{Теория \textperiodcentered{} законы цепей и энергия}
\hh{Закон Ома}
\fm{I = U/R} (или \fm{U = I R}). Сопротивление \vr{R} [Ом] — «трудность» протолкнуть ток.
У провода \fm{R = \rho L / A}: растёт с длиной, падает с сечением. Решатель считает цепь
честно, поэтому закон Ома выполняется во всех демо.

\hh{Законы Кирхгофа}
\textbf{Первый (токов).} В любом узле сумма токов равна нулю — заряд не копится: \fm{\sum I = 0}.
\textbf{Второй (напряжений).} Сумма падений напряжения вдоль замкнутого контура равна нулю.
На этих двух законах стоит \textbf{метод узловых потенциалов (MNA)} — им решатель и считает
схему.

\hh{Последовательно и параллельно}
Последовательно: ток общий, сопротивления складываются \fm{R = R_1 + R_2}, напряжение
делится. Параллельно: напряжение общее, складываются проводимости \fm{1/R = 1/R_1 + 1/R_2},
делится ток. Отсюда — \textbf{делитель напряжения} \fm{U_{\text{out}} = U\,R_2/(R_1{+}R_2)}
(работа~2), его \textbf{просадка под нагрузкой} (работа~3) и \textbf{мост Уитстона}: баланс
при \fm{R_1 R_4 = R_2 R_3} (работа~4).

\hh{Суперпозиция}
В линейной цепи отклик от нескольких источников равен сумме откликов от каждого по
отдельности (прочие источники мысленно «занулены»). Демо \textit{Superposition}.

\hh{Мощность и тепло}
Мощность \fm{P = U I} [Вт]. На сопротивлении она переходит в тепло (закон Джоуля–Ленца):
\fm{P = I^2 R = U^2/R}. Источник мощность отдаёт, резистор — рассеивает; слой \textit{power}
показывает знак и величину, \textit{heat} — нагрев. Температура считается простой моделью
теплообмена \fm{C_{th}\,dT/dt = P - (T-T_0)/R_{th}}, в стационаре \fm{T = T_0 + P\,R_{th}}.
\clearpage

\section*{Теория \textperiodcentered{} конденсатор, катушка, переходные процессы}
\hh{Конденсатор}
Запасает заряд в электрическом поле: \fm{q = C U}, энергия \fm{W = \tfrac12 C U^2}. Ток — это
скорость изменения заряда: \fm{i = C\,dU/dt}. Поэтому напряжение на конденсаторе не может
прыгнуть мгновенно. В установившемся постоянном токе ток через него равен нулю — «разрыв».

\hh{Зарядка RC}
Через резистор конденсатор заряжается по экспоненте: \fm{U_c(t) = U(1 - e^{-t/\tau})},
постоянная времени \fm{\tau = R C}. За \vr{\tau} — до \ensuremath{\approx}63\,\%, за 5\vr{\tau} —
почти полностью (работа~5).

\hh{Катушка индуктивности}
Запасает энергию в магнитном поле: потокосцепление \fm{\Psi = L I}, энергия
\fm{W = \tfrac12 L I^2}. Напряжение — скорость изменения тока: \fm{u = L\,dI/dt}. Поэтому ток
через катушку не прыгает мгновенно («инерция тока»). В постоянном токе — «провод» (КЗ).

\hh{Нарастание RL}
\fm{I(t) = (U/R)(1 - e^{-t/\tau})}, \fm{\tau = L/R} (работа~6).

\hh{Колебательный контур RLC}
Конденсатор и катушка обмениваются энергией (поле \ensuremath{\leftrightarrow} ток) — возникают
колебания с собственной частотой \fm{f_0 = 1/(2\pi\sqrt{LC})}. Резистор их гасит: чем он меньше,
тем дольше контур «звенит» (выше добротность \vr{Q}) (работа~7).

\hh{Как это считает симулятор}
Между моментами времени реактивные элементы заменяются эквивалентной схемой (companion-модель:
проводимость~+~источник, метод обратного Эйлера/трапеций), и на каждом шаге \vr{dt} снова решается
обычная линейная цепь. Постоянный ток — это предел при больших временах.
\clearpage

\section*{Теория \textperiodcentered{} переменный ток и диод}
\hh{Переменное напряжение}
Источник переменного тока даёт синусоиду \fm{v(t) = A\sin(2\pi f t + \varphi)}: амплитуда
\vr{A}, частота \vr{f} [Гц], фаза \vr{\varphi}. На реактивных элементах ток и напряжение
сдвинуты по фазе, а их «сопротивление» (реактивное) зависит от частоты.

\hh{Диод}
Пропускает ток практически в одну сторону. В прямом включении открывается выше порога
\fm{V_f \approx 0.7\ \text{В}} и дальше держит его примерно постоянным; в обратном — почти не
проводит (работа~8). Решатель моделирует диод кусочно-линейно (открыт/закрыт), подбирая
состояние итерацией.

\hh{Выпрямление}
Диод срезает отрицательную полуволну — на нагрузке остаются только положительные «горбы»
(полуволновой выпрямитель, работа~9). Конденсатор сглаживает их в почти постоянное напряжение
(пиковый детектор, RC-сглаживание).

\hh{Частотные фильтры}
RC-цепь под переменным сигналом — фильтр. Нижних частот пропускает медленные сигналы и режет
быстрые; частота среза \fm{f_c = 1/(2\pi R C)} (работа~10). Контур RLC выделяет узкую полосу
около резонанса \fm{f_r = 1/(2\pi\sqrt{LC})} (полосовой фильтр).
\clearpage

\section*{Теория \textperiodcentered{} магнитное поле тока}
\hh{Поле прямого провода}
Вокруг провода с током есть магнитное поле, охватывающее его кольцами. Для длинного прямого
провода \fm{B = \mu_0 I/(2\pi r)} — спадает с расстоянием \vr{r}. Направление — по правилу
правой руки (буравчика): обхвати провод, большой палец по току — пальцы покажут поле. Для
\textit{конечного} отрезка симулятор считает точную формулу Био–Савара
\fm{B = \dfrac{\mu_0 I}{4\pi r}(\sin\theta_1 + \sin\theta_2)}. Слой \textit{magnetic} рисует поле
с признаком «из/в плоскость».

\hh{Соленоид (катушка)}
Витки складывают свои поля, и внутри длинной катушки поле почти однородно:
\fm{B = \mu_0 n I = \mu_0 N I / l} (\vr{n} — витков на метр, \vr{l} — длина). Снаружи поле слабое
и похоже на поле магнитного диполя (спад \ensuremath{\sim}\vr{1/r^3}). Так катушка превращает ток
в магнит.

\hh{Что в модели НЕ входит}
Симулятор честно считает магнитное поле \emph{токов} (провод, соленоид), но не моделирует
электромагнитную индукцию как процесс, магнитные материалы и ЭМ-волны — катушка фигурирует лишь
как элемент цепи с инерцией тока.
\clearpage

\section*{Теория \textperiodcentered{} механическая и гидравлическая аналогии}
\lead{Одна и та же математика описывает разные системы — поэтому абстрактный ток можно
«пощупать». В симуляторе мощность и энергия аналогий совпадают с электрическими точно.}

\hh{Механическая аналогия (цепь-цепочка)}
Замкнутая цепь из звеньев на двух звёздочках (в узлах). Скорость звеньев = ток, натяжение =
потенциал. Резистор — тормоз (трение в тепло), конденсатор — пружина (сжатие = заряд), катушка —
маховик (инерция вращения = инерция тока).
\medskip
\begin{tabular}{@{}>{\bfseries}l l l@{}}
\toprule
Электрика & Механика & Связь \\
\midrule
ток \vr{I} & скорость цепи & \vr{v \propto I} \\
потенциал \vr{\varphi} & натяжение & \vr{T \propto \varphi} \\
резистор \vr{R} & тормоз (трение) & рассеяние \vr{= I^2R} \\
конденсатор \vr{C} & пружина & \vr{W=\tfrac12 C U^2 = \tfrac12 k x^2} \\
катушка \vr{L} & маховик & \vr{W=\tfrac12 L I^2 = \tfrac12 J \omega^2} \\
диод & храповик & пропуск в одну сторону \\
источник & кривошип (динамо) & задаёт натяжение \\
земля & неподвижный якорь & \vr{\varphi_{\text{ref}}=0} \\
\bottomrule
\end{tabular}

\hh{Гидравлическая аналогия}
Вода в трубах. Поток \vr{Q} = ток, давление = напряжение.
\medskip
\begin{tabular}{@{}>{\bfseries}l l l@{}}
\toprule
Электрика & Гидравлика & Связь \\
\midrule
ток \vr{I} & поток \vr{Q} & \vr{Q \propto I} \\
напряжение \vr{U} & перепад давления \vr{\Delta P} & \vr{\Delta P \propto U} \\
резистор \vr{R} & узкая труба & \vr{Q = \Delta P / R} \\
конденсатор \vr{C} & упругий бак & уровень \vr{\propto C\,U_c} \\
катушка \vr{L} & турбина / тяжёлая труба & инерция потока \\
диод & обратный клапан & поток в одну сторону \\
источник & насос & держит \vr{\Delta P} \\
земля & открытый резервуар & \vr{P_{\text{ref}}=0} \\
\bottomrule
\end{tabular}
\clearpage

% ================= РАБОТА 1 =================
\lab{1}{Закон Ома}
\lead{Связь тока, напряжения и сопротивления на одном резисторе.}
\hh{Цель}
Убедиться, что ток через резистор прямо пропорционален напряжению и обратно — сопротивлению.
\fig{lab01_ohm.png}{Источник 5~В и резистор 1~к\ensuremath{\Omega}. Стрелки — направление тока, цвет провода — потенциал.}
\hh{Теория}
\fm{I = V / R}\quad\fm{V = I\cdot R}\quad\fm{P = V\cdot I = I^2\cdot R}
\begin{predict}\textbf{Предскажи.} При \vr{V}\,=\,5~В и \vr{R}\,=\,1~к\ensuremath{\Omega} — чему равен ток? А если удвоить \vr{R} до 2~к\ensuremath{\Omega}?\end{predict}
\begin{observe}\textbf{Наблюдай.} Открой демо «source + resistor», сними \vr{I} с инспектора. Поменяй \vr{R} в редакторе элемента.\end{observe}
\begin{explain}\textbf{Объясни.} \vr{I} = 5/1000 = 5~мА; при 2~к\ensuremath{\Omega} — 2.5~мА (вдвое меньше).
\misc{Заблуждение:} «источник даёт фиксированный ток». Нет — источник держит \textit{напряжение}, а ток определяет нагрузка.\end{explain}
\clearpage

% ================= РАБОТА 2 =================
\lab{2}{Делитель напряжения}
\lead{Два последовательных резистора делят напряжение пропорционально сопротивлениям.}
\hh{Цель}
Понять, что в последовательной цепи \textbf{ток общий}, а напряжение \textbf{делится}.
\fig{lab02_divider.png}{Делитель: два резистора последовательно, выход — со средней точки.}
\hh{Теория}
\fm{V_{\text{out}} = V \cdot R_2/(R_1{+}R_2)}\quad ток один на всю ветвь: \fm{I = V/(R_1{+}R_2)}
\begin{predict}\textbf{Предскажи.} \vr{V}\,=\,5~В, \vr{R_1}\,=\,\vr{R_2}\,=\,1~к\ensuremath{\Omega} — сколько на выходе? А если \vr{R_2}\,=\,4~к\ensuremath{\Omega}?\end{predict}
\begin{observe}\textbf{Наблюдай.} Демо «resistor divider». Сними потенциал средней точки. Меняй \vr{R_2}.\end{observe}
\begin{explain}\textbf{Объясни.} Поровну → 2.5~В; при \vr{R_2}\,=\,4~к\ensuremath{\Omega} → $5\cdot4/5 = 4$~В.
\misc{Заблуждение:} «ток расходуется на первом резисторе». Нет — ток через оба одинаков; «расходуется» (падает) \textit{напряжение}.\end{explain}
\clearpage

% ================= РАБОТА 3 =================
\lab{3}{Эффект нагрузки}
\lead{Почему идеальный делитель «проседает» под реальной нагрузкой.}
\hh{Цель}
Увидеть, что подключение нагрузки параллельно плечу делителя снижает выходное напряжение.
\fig{lab03_loaded.png}{Делитель \ensuremath{R_1}/\ensuremath{R_2} с нагрузкой \ensuremath{R_L} параллельно \ensuremath{R_2}.}
\hh{Теория}
Нагрузка параллельна \vr{R_2}: \fm{R_2\parallel R_L = R_2\cdot R_L/(R_2{+}R_L)}, и эта величина \textbf{меньше} \vr{R_2} → выход падает.
\begin{predict}\textbf{Предскажи.} Без нагрузки \vr{V_{\text{out}}}\,\ensuremath{\approx}\,3.33~В (\vr{R_1}=1\,к, \vr{R_2}=2\,к). Что станет с \vr{V_{\text{out}}}, когда повесим \vr{R_L}\,=\,1~к\ensuremath{\Omega} — вырастет, упадёт, не изменится?\end{predict}
\begin{observe}\textbf{Наблюдай.} Демо «loaded voltage divider». Сравни выход с расчётом без нагрузки.\end{observe}
\begin{explain}\textbf{Объясни.} $R_2\parallel R_L = 2000\cdot1000/3000 \approx 667$~\ensuremath{\Omega} → \vr{V_{\text{out}}} $= 5\cdot667/(1000+667) \approx 2.0$~В. Просадка.
\misc{Заблуждение:} «делитель даёт фиксированный выход независимо от нагрузки».\end{explain}
\clearpage

% ================= РАБОТА 4 =================
\lab{4}{Мост Уитстона}
\lead{Точное измерение через баланс, а не абсолютный отсчёт.}
\hh{Цель}
Понять условие баланса моста и почему через «гальванометр» течёт ток при разбалансе.
\fig{lab04_wheatstone.png}{Четыре плеча \ensuremath{R_1..R_4} и мостовой резистор \ensuremath{R_g} между средними точками.}
\hh{Теория}
Баланс (ток через \vr{R_g} равен нулю): \fm{R_1\cdot R_4 = R_2\cdot R_3}. При нарушении — появляется ток через мост.
\begin{predict}\textbf{Предскажи.} \vr{R_1}=\vr{R_2}=\vr{R_3}=1~к\ensuremath{\Omega}, \vr{R_4}=2~к\ensuremath{\Omega}. Сбалансирован ли мост? Потечёт ли ток через \vr{R_g}?\end{predict}
\begin{observe}\textbf{Наблюдай.} Демо «Wheatstone bridge». Сними ток через \vr{R_g}. Сделай \vr{R_4}=1~к\ensuremath{\Omega} — что изменится?\end{observe}
\begin{explain}\textbf{Объясни.} $1\cdot2 \neq 1\cdot1$ → разбаланс → ток через \vr{R_g}. При \vr{R_4}=1~к\ensuremath{\Omega}: $1\cdot1 = 1\cdot1$ → баланс → ток нулевой.
\misc{Заблуждение:} «мост всегда измеряет напряжение». Он измеряет \textit{отношение} сопротивлений по нулю тока.\end{explain}
\clearpage

% ================= РАБОТА 5 =================
\lab{5}{RC \textperiodcentered{} постоянная времени}
\lead{Конденсатор заряжается не мгновенно, а по экспоненте.}
\hh{Цель}
Измерить постоянную времени \vr{\tau} и связать её с \vr{R} и \vr{C}.
\figtwo{lab05_rc.png}{Электрическая: источник → R → C.}{lab05_rc_mech.png}{Механика: конденсатор = пружина (заряд = сжатие).}
\hh{Теория}
\fm{V_c(t) = V\cdot(1 - e^{-t/\tau})}\quad\fm{\tau = R\cdot C}. За \vr{\tau} заряжается до \ensuremath{\approx}63\,\%, за 5\vr{\tau} — почти полностью.
\begin{predict}\textbf{Предскажи.} \vr{R}=1~к\ensuremath{\Omega}, \vr{C}=1~мФ. Чему равна \vr{\tau}? Сколько времени до \ensuremath{\approx}63\,\% от 5~В?\end{predict}
\begin{observe}\textbf{Наблюдай.} Демо «RC charging». Запусти, следи за \vr{V_c(t)} на осциллографе. В механическом виде смотри, как пружина плавно сжимается.\end{observe}
\begin{explain}\textbf{Объясни.} \vr{\tau} $= 1000\cdot0.001 = 1$~с; до 63\,\% (\ensuremath{\approx}3.16~В) — ровно за 1~с.
\misc{Заблуждение:} «конденсатор заряжается мгновенно» или «хранит напряжение, а не заряд». Хранит заряд \vr{Q = C\cdot V}.\end{explain}
\clearpage

% ================= РАБОТА 6 =================
\lab{6}{RL \textperiodcentered{} инерция тока}
\lead{Катушка сопротивляется \textit{изменению} тока — ток нарастает плавно.}
\hh{Цель}
Понять, что индуктивность задаёт инерцию току, и измерить \fm{\tau = L/R}.
\figtwo{lab06_rl.png}{Электрическая: источник → R → L (катушка — витки).}{lab06_rl_mech.png}{Механика: катушка = маховик (ток = момент импульса). Зубчатый обод сцеплен с цепью.}
\hh{Теория}
\fm{I(t) = (V/R)\cdot(1 - e^{-t/\tau})}\quad\fm{\tau = L/R}. В стационаре катушка = провод (КЗ).
\begin{predict}\textbf{Предскажи.} \vr{V}=5~В, \vr{R}=10~\ensuremath{\Omega}, \vr{L}=1~Гн. Конечный ток? \vr{\tau}? Мгновенно ли он установится?\end{predict}
\begin{observe}\textbf{Наблюдай.} Демо «RL current rise». Следи за током. В механике маховик-катушка раскручивается постепенно — это и есть инерция тока.\end{observe}
\begin{explain}\textbf{Объясни.} \vr{I_\infty} = 5/10 = 0.5~А; \vr{\tau} = 1/10 = 0.1~с. Ток нарастает, а не прыгает.
\misc{Заблуждение:} «ток в катушке появляется сразу». Маховик не раскрутить мгновенно.\end{explain}
\clearpage

% ================= РАБОТА 7 =================
\lab{7}{Колебательный контур RLC}
\lead{Конденсатор и катушка вместе дают колебания — энергия перетекает туда-обратно.}
\hh{Цель}
Увидеть затухающие колебания и связать их частоту с \vr{L} и \vr{C}.
\fig{lab07_rlc.png}{Последовательный RLC во время переходного процесса (\ensuremath{t \approx 1}~с).}
\hh{Теория}
Частота собственных колебаний: \fm{f_0 = 1/(2\pi\cdot\sqrt{L\cdot C})}. Резистор гасит колебания (затухание).
\begin{predict}\textbf{Предскажи.} Что будет с током после включения: монотонно установится или «зазвенит»? От чего зависит частота звона?\end{predict}
\begin{observe}\textbf{Наблюдай.} Демо «RLC series». Следи за током/напряжением на осциллографе — увидишь затухающую синусоиду. В механике пружина (C) и маховик (L) раскачивают цепь.\end{observe}
\begin{explain}\textbf{Объясни.} Энергия колеблется между полем конденсатора (\vr{\tfrac12 CV^2}) и движением маховика (\vr{\tfrac12 LI^2}); R забирает её в тепло → амплитуда падает.
\misc{Заблуждение:} «реактивные элементы просто замедляют». Вместе они \textit{колеблют}.\end{explain}
\clearpage

% ================= РАБОТА 8 =================
\lab{8}{Диод}
\lead{Элемент, пропускающий ток только в одну сторону — с порогом.}
\hh{Цель}
Понять одностороннюю проводимость и прямое падение \ensuremath{\approx}0.7~В.
\fig{lab08_diode.png}{Диод последовательно с резистором.}
\hh{Теория}
Прямое смещение (анод «+»): проводит, падение \fm{V_f \approx 0.7\ \text{В}}. Обратное: не проводит (\ensuremath{\approx} обрыв).
\begin{predict}\textbf{Предскажи.} Источник 5~В через диод на резистор 1~к\ensuremath{\Omega}. Какое напряжение упадёт на резисторе — 5~В или меньше? Если перевернуть диод — потечёт ли ток?\end{predict}
\begin{observe}\textbf{Наблюдай.} Демо «diode + resistor». Сними падение на диоде и на резисторе. Переверни диод (поменяй узлы).\end{observe}
\begin{explain}\textbf{Объясни.} На резисторе \ensuremath{\approx}5 - 0.7 = 4.3~В (\ensuremath{\approx}4.3~мА). Перевёрнутый диод — тока почти нет.
\misc{Заблуждение:} «диод — идеальный ключ без потерь». Есть порог \ensuremath{\approx}0.7~В.\end{explain}
\clearpage

% ================= РАБОТА 9 =================
\lab{9}{Полуволновой выпрямитель}
\lead{Как из переменного напряжения получить пульсирующее постоянное.}
\hh{Цель}
Увидеть, что диод срезает отрицательную полуволну, оставляя ток в одну сторону.
\figtwo{lab09_ac_pos.png}{Положительная полуволна: диод проводит, ток течёт.}{lab09_ac_neg.png}{Отрицательная полуволна: диод заперт, ток не течёт.}
\hh{Теория}
Источник \fm{v(t) = A\cdot\sin(2\pi\cdot f\cdot t)}. Диод пропускает только \vr{v}\,>\,\vr{V_f} → на нагрузке — однополярные импульсы. Конденсатор сглаживает.
\begin{predict}\textbf{Предскажи.} На входе синус \ensuremath{\pm}5~В, 2~Гц. Каким будет напряжение на нагрузке — синус, модуль синуса, или только положительные «горбы»?\end{predict}
\begin{observe}\textbf{Наблюдай.} Демо «AC half-wave rectifier». Сравни два момента времени (положительная и отрицательная фаза) — на отрицательной диод заперт.\end{observe}
\begin{explain}\textbf{Объясни.} Выход — только положительные полупериоды (минус срезан). Среднее\,>\,0 → есть постоянная составляющая.
\misc{Заблуждение:} «переменный ток нельзя превратить в постоянный без сложной электроники».\end{explain}
\clearpage

% ================= РАБОТА 10 =================
\lab{10}{RC-фильтр НЧ}
\lead{Та же RC-цепь, но под переменным сигналом — становится частотным фильтром.}
\hh{Цель}
Понять частоту среза и почему высокие частоты ослабляются сильнее.
\fig{lab10_lowpass.png}{RC-фильтр нижних частот под синусоидальным источником.}
\hh{Теория}
Частота среза: \fm{f_c = 1/(2\pi\cdot R\cdot C)}. Ниже \vr{f_c} сигнал проходит, выше — слабеет (на \vr{f_c} выход \ensuremath{\approx}0.71 от входа).
\begin{predict}\textbf{Предскажи.} \vr{R}=1~к\ensuremath{\Omega}, \vr{C}=100~мкФ. Чему равна \vr{f_c}? При входе 2~Гц выход будет почти равен входу или заметно ослаблен?\end{predict}
\begin{observe}\textbf{Наблюдай.} Демо «RC low-pass filter (AC)». Сравни амплитуду на C (выход) с амплитудой источника.\end{observe}
\begin{explain}\textbf{Объясни.} \vr{f_c} $= 1/(2\pi\cdot1000\cdot100\cdot10^{-6}) \approx 1.59$~Гц. Вход 2~Гц чуть выше среза → выход заметно ослаблен ($\sim$0.62$\times$).
\misc{Заблуждение:} «конденсатор просто запасает заряд». На переменном токе он — частотно-зависимое сопротивление.\end{explain}
\clearpage

% ================= ПРИЛОЖЕНИЕ =================
\section*{Приложение \textperiodcentered{} как воспроизвести иллюстрации}
Все схемы этого практикума сняты прямо из симулятора, без скриншотов ОС, режимом
оффскрин-съёмки. Пример:

\texttt{current-lab -{}-capture -{}-demo RcCapacitor -{}-view electrical -{}-out rc.png}\\
\texttt{current-lab -{}-capture -{}-demo RlcSeries -{}-view mechanical -{}-time 1.0 -{}-out rlc.png}

\medskip
\begin{tabular}{@{}>{\bfseries}l p{0.74\linewidth}@{}}
\toprule
Флаг & Значение \\
\midrule
\texttt{-{}-demo} & имя демо-схемы (enum-имя или подпись «Demo: …») \\
\texttt{-{}-view} & electrical \textperiodcentered{} mechanical \textperiodcentered{} hydraulic \\
\texttt{-{}-time} & момент СИМУЛИРОВАННОГО времени, с (по умолчанию — до стационара) \\
\texttt{-{}-out}  & путь к PNG \\
\bottomrule
\end{tabular}

\vspace{24pt}
{\color{muted}\small Electricity Lab \textperiodcentered{} MNA-решатель \textperiodcentered{} 2026-06-15.
Метод «предскажи → наблюдай → объясни» основан на predict-observe-explain (POE).}

\end{document}
